\section{The 41\% of work that is thrown away \hfill \live{}}

\textbf{The observation.} We store the minimum possible number of states. We do \emph{not} expand
the minimum possible number. The difference between states explored and states stored is exactly
the number of stored zones that were later retired --- states we did the full work of expanding,
only for a better zone to arrive afterwards and make them redundant.

\begin{center}
\begin{tabular}{@{}lrr@{}}
\toprule
model & wasted expansions & share \\
\midrule
\texttt{fischer10} & $186\,600$ & $41.7\%$ \\
\texttt{twolock10} & --- & $\approx 41\%$ \\
\texttt{phil12}    & $224\,030$ & $27.9\%$ \\
\texttt{arbiter14}, \texttt{viking12}, \texttt{csma} & $0$ & $0\%$ \\
\bottomrule
\end{tabular}
\end{center}

\paragraph{Why this one is different from everything else here.} It is not benchmark-dependent, and
that is not a hope, it is measured. The delay between a state being expanded and being retired is a
\emph{constant of the model}: exactly $4$ on \texttt{fischer} at both $8$ and $10$ processes,
exactly $6$ on \texttt{twolock}, multiples of $4$ on \texttt{phil}. And the constant is
\emph{syntactic} --- it equals the length of the shortest cycle in a single process template that
returns to the same discrete key. A clock-reset side condition then predicts whether the tax exists
at all, and it predicted presence or absence correctly on $6$ of $6$ families \emph{before} running
the search.

So the precondition is checkable on any input model by looking at the templates. That is the
generality bar the project set, and this is the only lead that clears it by construction.

\paragraph{The catch, and it is a real one.} Simply holding every state back for $g$ steps is
provably useless --- the arithmetic works out so that a uniform delay costs at least as much as it
saves. This is not a small detail: it is exactly the effect Behrmann et al.\ measured in 2000 when
they distributed \uppaal{} and watched the state count \emph{rise} from $45\,001$ on one node to
$110\,634$ on seventeen. Any deferral scheme has to be selective, and the selection has to be
driven by the syntactic constant rather than by a tuned parameter.

\paragraph{The one thing that could kill it.} If the zone that retires a state is a
\emph{descendant} of that state, then the state had to be expanded to produce its own killer, and
no amount of waiting helps. A probe answered this: of $40$ sampled cases, \textbf{zero} required
the state's own expansion, and an oracle-order test attributes $86$--$89\%$ of the waste to
ordering. So the idea survives its own worst case.

\paragraph{Where parallelism comes in.} Deferral is a bet: pay latency now to avoid a subtree
later. A sequential search cannot afford the bet, because delaying its only state stalls it
completely. With many workers there is always other work, so the bet is free. This is the one place
in the whole document where extra cores buy something a sequential tool cannot have at any price.
